\chapter{Introduction}\label{ch:introduction}

% Introduce the topic of the thesis, including research questions.

\section{Motivation and research questions}

% TODO: Give a bit of background on lakehouses, their design and the issues they aim to solve. Explain how each of the three lakehouses works and then transition into the research questions.

Many modern data storage architectures rely on so-called \textit{data lakes}, which consist of cloud object stores, e.g. Amazon S3\footnote{\url{https://aws.amazon.com/s3/}}, Azure Blob Storage\footnote{\url{https://azure.microsoft.com/en-us/products/storage/blobs/}}, for storing large amounts of (un)structured and semi-structured data in a centralized place \cite{WhatDataLake}. Part of the data can then be copied into data warehouses through a process called "extract, transform, load" (ETL), so that it can be used for further downstream analytics tasks \cite{michael_armbrust_lakehouse_2021}. Although this two-tier architecture has become widespread among many enterprises, it has a few drawbacks. ETL pipelines used to load data into warehouses can be time-consuming and error-prone, causing the data in the warehouse to be outdated and sometimes inconsistent \cite{schneider_lakehouse_2024}. Additionally, maintaining the two systems is expensive and redundant, since part of the data is stored in both the data lake and the warehouse \cite{michael_armbrust_lakehouse_2021}. Finally, the flexibility of storing various file formats in the same place comes with the cost of managing different file versions and metadata, which with time can transform the data lake into a "data swamp" \cite{avrilayshaDeltaLakeVs}.
 
Following the limitations of data lakes and the two-tier architecture, \citeauthor{michael_armbrust_lakehouse_2021} \cite{michael_armbrust_lakehouse_2021} proposed a new data storage architecture called \textit{data lakehouse}. A lakehouse, also referred to as an open table format, builds on top of data lakes by structuring data into tables that are stored in immutable columnar formats, e.g. Apache Parquet\footnote{\url{https://parquet.apache.org/}}. These tables are managed with centralized metadata files that keep track of the current schema and the different versions of the table, all while ensuring that updates to tables occur in ACID \cite{haerderPrinciplesTransactionorientedDatabase1983} transactions \cite{michael_armbrust_lakehouse_2021}. Furthermore, the use of open formats for storage gives users control over their data, in contrast to proprietary formats of data warehouses \cite{michael_armbrust_lakehouse_2021, schneider_lakehouse_2024, avrilayshaDeltaLakeVs}. These characteristics allow analytical tasks to be performed directly on the lakehouse tables, thus eliminating the need of having a separate data warehouse, thus eliminating the need for a two-tier architecture.

Although open table formats have gained popularity since their introduction, some aspects of lakehouses have led to important questions in recent works. For instance, \citeauthor{paras_jain_analyzing_2023} \cite{paras_jain_analyzing_2023} suggest that performing high frequent insertions into lakehouse tables can be challenging, because every insertion creates a new data file and requires updating the metadata files as well. Further, inserts in Delta Lake, a popular lakehouse format, can have up to hundreds of milliseconds of latency due to its reliance on object stores, which limits the throughput of streaming workloads, as mentioned in \cite{armbrust_delta_2020}. By contrast, the recent introduction of DuckLake \cite{mark_raasveldt_ducklake_2025} brought significant changes on how to handle write transactions to tables, while still ensuring ACID transactions. This is all thanks to the use of a traditional database management system (DBMS), e.g. PostgreSQL\footnote{\url{https://www.postgresql.org/}}, for handling table metadata. Interestingly, no work has been done yet to evaluate the throughput of open table formats when performing data stream insertions. Therefore, this thesis will fill that gap, by focusing on the following research question:

% Consequently, this raises the question of just how well-suited current lakehouse systems really are at meeting data stream processing (DSP) requirements. For example, the authors in \cite{paras_jain_analyzing_2023} conducted an evaluation of three different lakehouse systems on the TPC-DS benchmark \cite{TPCBENCHMARKDS2024} and suggested to further investigate the ability of lakehouses at handling workflows with a high number of concurrent write queries per second. They believe that this would be a challenging task for lakehouse systems due to the relatively high latency of communicating with the object store for updating the metadata on every write operation \cite{paras_jain_analyzing_2023}. Additionally, the authors in \cite{armbrust_delta_2020} mention that inserts in Delta can have up to hundreds of milliseconds of latency due to its reliance on object stores' implementations of put-if-absent operations for guaranteeing atomicity of transaction logs. Despite this, \citeauthor{armbrust_delta_2020} claim that Delta suffices in providing fast enough execution for all its applications \cite{armbrust_delta_2020}. Meanwhile, DuckLake's design might result in lower latency writes, since all the metadata is written to traditional DBMSs, which are more efficient at handling multiple concurrent transactions.

\begin{center}
    "\textit{How can we evaluate the throughput of data streams in lakehouse systems?}"
\end{center}

% Alternative
% \begin{center}
%     "\textit{How can we systematically evaluate streaming write performance in lakehouse systems?}"
% \end{center}

In addition, the following sub-questions will be answered 
\begin{enumerate}
    \item What are the current state-of-the-art benchmarks for lakehouses?
    % \begin{enumerate}
    %     \item What metrics are relevant for evaluating the performance of data streaming writes?
    % \end{enumerate}
    % \item What are the state-of-the-art datasets used for benchmarking streaming systems?
    \item What are common patterns used in benchmarks for stream processing systems?
    \item How can we ensure our benchmark is actually fair?
    \item What is the difference in performance of current lakehouse systems when evaluated with our benchmark?
    \begin{enumerate}
        \item What is the performance gain of the "data inlining" feature of DuckLake in comparison to other lakehouse formats?
    \end{enumerate}
\end{enumerate}

\section{Contributions}


\section{Thesis structure}
The rest of this thesis is structured as follows: Chapter \ref{ch:background} goes over the relevant background for this thesis, such as the different lakehouse technologies, how they are implemented and relevant features for this research. Chapter \ref{ch:related-work} talks about the existing work on benchmarking lakehouses and the different metrics employed by these benchmarks and the systems that they evaluate. Chapter \ref{ch:methodology} describes the methodology employed to develop our own benchmark and the reasoning behind those decisions. We also describe our experimental setup here. Chapter \ref{ch:results} present the results of our experiments, which will be discussed in Chapter \ref{ch:discussion}. Finally, Chapter \ref{ch:conclusion} will conclude this thesis by looking back at the proposed research questions, while proposing possible directions for future work.
